\documentclass[a4paper, UTF-8]{article}
\usepackage[margin=1in]{geometry}
\usepackage{algorithm}
\usepackage{algpseudocode}
\usepackage{amsmath}
\usepackage{amsfonts}
\usepackage{amsthm}
\usepackage{tikz}
\usetikzlibrary{arrows.meta, calc}

\usepackage[utf8]{inputenc} % allow utf-8 input
\usepackage[T1]{fontenc}    % use 8-bit T1 fonts
\usepackage{hyperref}       % hyperlinks
\usepackage{url}            % simple URL typesetting
\usepackage{booktabs}       % professional-quality tables
\usepackage{amsfonts}       % blackboard math symbols
\usepackage{microtype}      % microtypography
\usepackage{xcolor}         % colors

\theoremstyle{plain}
\newtheorem{theorem}{\indent Theorem}[section]
\newtheorem{lemma}[theorem]{\indent Lemma}
\newtheorem{assumption}{\indent Assumption}[section]
\newtheorem{note}[theorem]{\indent Notation}
\newtheorem{proposition}[theorem]{\indent Proposition}
\newtheorem{corollary}[theorem]{\indent Corollary}
\newtheorem{definition}{\indent Definition}[section]
\newtheorem{example}{\indent Example}[section]
\newtheorem{remark}{\indent Remark}[section]
\newenvironment{solution}{\begin{proof}[\indent\bf Solution]}{\end{proof}}
\renewcommand{\proofname}{\indent\bf Proof}

\begin{document}

\title{\textbf{Assignment 12}}

\author{Yanqiao Chen\\ 
  Department of Computer Science and Engineering\\
  Southern University of Science and Technology\\
  \texttt{12412115@mail.sustech.edu.cn}
  }


\maketitle

\section{Page 171, Problem 54}

Given that $X\sim N(\mu_X, \sigma_X^2)$, $Y\sim N(\mu_Y, \sigma_Y^2)$, $Z \sim N(\mu_Z, \sigma_Z^2)$, where they are uncorrelated, for
\begin{align*}
U = Z + X \\ 
V = Z + Y
\end{align*}
we have
\begin{align*}
\text{Cov}(U, V) & = \text{Cov}(Z + X, Z + Y) \\
& = \text{Cov}(Z, Z) + \text{Cov}(Z, Y) + \text{Cov}(X, Z) + \text{Cov}(X, Y) \\
& = \sigma_Z^2
\end{align*}
and
\begin{align*}
\rho_{UV} & = \frac{\text{Cov}(U, V)}{\sqrt{\text{Var}(U)\text{Var}(V)}} \\
& = \frac{\sigma_Z^2}{\sqrt{(\sigma_X^2 + \sigma_Z^2)(\sigma_Y^2 + \sigma_Z^2)}}
\end{align*}

\section{Page 171, Problem 60} 

Given that the density function have the symmetric property $f(y) = f(-y)$ then $E(Y) = 0$ and $X = SY$, where $S$ is a random variable independent variable taking on the value $+1$ and $-1$ with probability $\frac{1}{2}$, we have
\begin{align*}
\text{Cov}(X, Y) & = E(XY) - E(X)E(Y) \\
& = E(SY^2) - E(S)E(Y)E(Y) \\
& = E(S)E(Y^2) - E(S)E(Y)E(Y) \\
& = 0
\end{align*}

We show that $X$ and $Y$ are not independent. 
Since $X=SY$,
\begin{align*}
P(X=Y) & = P(S=1) + P(S=-1, Y=0) \\
& = \tfrac{1}{2} + 0 = \tfrac{1}{2},
\end{align*}
because $Y$ has a density and thus $P(Y=0)=0$. If $X$ and $Y$ were independent and each had a density, then $P(X=Y)=0$, a contradiction. Hence $X$ and $Y$ are not independent.

\section{Extra Question} 

\subsection{1)} 

\begin{align*}
E(X) & = \int_{-\infty}^{\infty} \frac{x}{2} e^{-|x|} \mathrm{d}x \\
& = \int_{0}^{\infty} \frac{x}{2} e^{-x} \mathrm{d}x + \int_{-\infty}^{0} \frac{x}{2} e^{x} \mathrm{d}x \\
& = 0
\end{align*}

\begin{align*}
\text{Var}(X) & = E(X^2) - E(X)^2 \\
& = \int_{-\infty}^{\infty} \frac{x^2}{2} e^{-|x|} \mathrm{d}x \\
& = \int_{0}^{\infty} \frac{x^2}{2} e^{-x} \mathrm{d}x + \int_{-\infty}^{0} \frac{x^2}{2} e^{x} \mathrm{d}x \\
& = 2
\end{align*}

\subsection{2)} 

No, $X$ and $|X|$ are not independent. If we know the value of $|X|$, the value of $X$ is completely restricted to $\pm |X|$. For instance, 
$$ P(X \ge 1 \mid |X| < 1) = 0 $$
but $P(X \ge 1) > 0$. Since the conditional probability is not equal to the marginal probability, they are dependent.

\subsection{3)} 

They are uncorrelated. To check for correlation, we compute the covariance:
\begin{align*}
\text{Cov}(X, |X|) & = E(X|X|) - E(X)E(|X|)
\end{align*}
We already computed that $E(X) = 0$. Furthermore,
\begin{align*}
E(X|X|) & = \int_{-\infty}^{\infty} x|x| \frac{1}{2} e^{-|x|} \mathrm{d}x
\end{align*}
Notice that $g(x) = x|x|$ is an odd function, and the density $f(x) = \frac{1}{2} e^{-|x|}$ is an even function. Their product is an odd function. Integrating an odd function over the symmetric interval $(-\infty, \infty)$ gives $0$. Hence, $E(X|X|) = 0$. 

Therefore,
\begin{align*}
\text{Cov}(X, |X|) & = 0 - 0 \cdot E(|X|) = 0 
\end{align*}
Since the covariance is zero, $X$ and $|X|$ are uncorrelated (not correlated).

\section{Supplementary Question 1.1} 

$$
f_X(x) = \int_{0}^{2} \frac{x+y}{8} \mathrm{d}y = \frac{x}{4} + \frac{1}{4}
$$

$$
f_Y(y) = \int_{0}^{2} \frac{x+y}{8} \mathrm{d}x = \frac{y}{4} + \frac{1}{4}
$$

$$
E(X) = \int_{0}^{2} x\left(\frac{x}{4} + \frac{1}{4}\right) \mathrm{d}x = \frac{7}{6}
$$

$$
E(Y) = \int_{0}^{2} y\left(\frac{y}{4} + \frac{1}{4}\right) \mathrm{d}y = \frac{7}{6} 
$$

$$
\text{Cov}(X, Y) = E(XY) - E(X)E(Y) = \int_{0}^{2}\int_{0}^{2} xy \frac{x+y}{8} \mathrm{d}x\mathrm{d}y - \frac{7}{6}\cdot\frac{7}{6} = \frac{4}{3} - \frac{49}{36} = -\frac{1}{36}
$$

\begin{align*}
\rho_{XY} &= \frac{\text{Cov}(X, Y)}{\sqrt{\text{Var}(X)\text{Var}(Y)}} \\ 
&= \frac{-\frac{1}{36}}{\sqrt{\left(\int_{0}^{2} x^2\left(\frac{x}{4} + \frac{1}{4}\right) \mathrm{d}x - \left(\frac{7}{6}\right)^2\right)\left(\int_{0}^{2} y^2\left(\frac{y}{4} + \frac{1}{4}\right) \mathrm{d}y - \left(\frac{7}{6}\right)^2\right)}} \\ 
&= \frac{-\frac{1}{36}}{\sqrt{\left(\frac{5}{3} - \frac{49}{36}\right)^2}} \\ 
& = -\frac{1}{11}
\end{align*}

\begin{align*}
D(X + Y) &= D(X) + D(Y) + 2\text{Cov}(X, Y) \\ 
&= \frac{11}{36} + \frac{11}{36} + 2\cdot\left(-\frac{1}{36}\right)\\ &= \frac{5}{9}
\end{align*}

\section{Supplementary Question 1.2}

Given that $X, Y \sim N(\mu, sigma)$ and $X, Y$ are independent. Let

$$
Z = \alpha X + \beta Y, W = \alpha X - \beta Y
$$

then 

\begin{align*}
\text{Cov}(Z, W) &= E(ZW) - E(Z)E(W) \\ 
&= E((\alpha X + \beta Y)(\alpha X - \beta Y)) - E(\alpha X + \beta Y)E(\alpha X - \beta Y) \\ 
&= E(\alpha^2 X^2 - \beta^2 Y^2) - E(\alpha X + \beta Y)E(\alpha X - \beta Y) \\
&= \alpha^2 E(X^2) - \beta^2 E(Y^2) - (\alpha E(X) + \beta E(Y))(\alpha E(X) - \beta E(Y)) \\
&= \alpha^2 (\sigma^2 + \mu^2) - \beta^2 (\sigma^2 + \mu^2) - (\alpha \mu + \beta \mu)(\alpha \mu - \beta \mu) \\
&= (\alpha^2 - \beta^2)(\sigma^2 + \mu^2) - (\alpha^2 - \beta^2)\mu^2 \\
&= (\alpha^2 - \beta^2)\sigma^2
\end{align*}

\begin{align*} 
\rho_{ZW} &= \frac{\text{Cov}(Z, W)}{\sqrt{\text{Var}(Z)\text{Var}(W)}} \\
&= \frac{(\alpha^2 - \beta^2)\sigma^2}{\sqrt{(\alpha^2 \sigma^2 + \beta^2 \sigma^2)(\alpha^2 \sigma^2 + \beta^2 \sigma^2)}} \\
&= \frac{(\alpha^2 - \beta^2)\sigma^2}{\sqrt{(\alpha^2 + \beta^2)^2 \sigma^4}} \\
&= \frac{\alpha^2 - \beta^2}{\alpha^2 + \beta^2}
\end{align*}

\section{Question 1} 

Let 
$$
Z = X - Y
$$

Thus $E(Z) = 2 - 2 = 0$, and
\begin{align*}
\text{Var}(Z) & = \text{Var}(X) + \text{Var}(Y) - 2\rho \sqrt{\text{Var}(X)\text{Var}(Y)} \\
& = 1 + 4 - 2\cdot\frac{1}{2}\cdot \sqrt{1 \cdot 4} \\ 
& = 3
\end{align*}

Thus 
$$
P(|X - Y| \geq 6) = P(|Z - E(Z)| \geq 6) \leq \frac{\text{Var}(Z)}{6^2} = \frac{1}{12}
$$

\section{Question 2} 

Let $X_1, X_2, \ldots, X_n$ be the number of dice rolled, respectively. The arithmetic mean of the number of dice rolled is $\bar{X} = \frac{1}{n}\sum_{i=1}^{n} X_i$. We have $E(X_i) = 3.5$. By the weak law of large numbers, the mean $\bar{X}$ converges in probability to $E(X_i) = 3.5$. 

\section{Question 3} 

Let the random variable $X \sim EXP(100)$, picking 16 samples randomly from the components. Thus, 
\begin{align*}
P\left(\sum_{i=1}^{16} (X_i) > 1920\right) &= P\left(\left(\frac{\sum_{i=1}^{16} (X_i) - 16\cdot E(X)}{\sqrt{16\cdot \text{Var}(X)}} > \frac{1920 - 16\cdot E(X)}{\sqrt{16\cdot \text{Var}(X)}}\right)\right) \\
& = P\left(\frac{\sum_{i=1}^{16} (X_i) - 16\cdot E(X)}{\sqrt{16\cdot \text{Var}(X)}} > \frac{1920 - 1600}{400}\right) \\
& = P\left(\frac{\sum_{i=1}^{16} (X_i) - 16\cdot E(X)}{\sqrt{16\cdot \text{Var}(X)}} > 0.8\right) \\
& = 1 - P\left(\frac{\sum_{i=1}^{16} (X_i) - 16\cdot E(X)}{\sqrt{16\cdot \text{Var}(X)}} \leq 0.8\right) \\
& = 1 - \Phi(0.8) \\
& \approx 0.2119
\end{align*}

\section{Question 4} 

Let $S$ be the number of death in a year, $E(S) = 10000 \times 0.017 = 170$. $\text{Var}(S) = 10000 \times 0.017 \times 0.983 = 167.11$. 

Thus, 
\begin{align*}
P(S > 200) & = P\left(\frac{S - E(S)}{\sqrt{\text{Var}(S)}} > \frac{200 - 170}{\sqrt{167.11}}\right) \\
& = P\left(\frac{S - E(S)}{\sqrt{\text{Var}(S)}} > 2.32\right) \\
& = 1 - P\left(\frac{S - E(S)}{\sqrt{\text{Var}(S)}} \leq 2.32\right) \\
& = 1 - \Phi(2.32) \\
& \approx 0.0102
\end{align*}

\section{Supplementary Question 2.1}

For $X_k$, $P(X_k = \sqrt{\ln k}) = \frac{1}{2}$, and $P(X_k = -\sqrt{\ln k}) = \frac{1}{2}$. Thus, $E(X_k) = 0$ and $\text{Var}(X_k) = \ln k$. We prove that it follows the law of large numbers by checking the sufficient condition.
\begin{align*}
\lim_{n\to\infty} \frac{1}{n^2}\sum_{k=1}^{n} \text{Var}(X_k) & = \lim_{n\to\infty} \frac{1}{n^2}\sum_{k=1}^{n} \ln k \\
& \leq \lim_{n\to\infty} \frac{1}{n^2} \cdot n \ln n \\
& = \lim_{n\to\infty} \frac{\ln n}{n} \\
& = 0
\end{align*}
Thus, by the sufficient condition of the weak law of large numbers, we have $\bar{X} = \frac{1}{n}\sum_{k=1}^{n} X_k$ converges in probability to $E(X_k) = 0$.

\section{Supplementary Question 2.2} 

For $E(X_n) = p_n$, $\text{Var}(X_n) = p_n(1-p_n)$. We check the sufficient condition of the weak law of large numbers:

\begin{align*} 
\lim_{n \to \infty} \frac{1}{n^2} \text{Var}(\sum_{i=1}^{n} X_i) & = \lim_{n \to \infty} \frac{1}{n^2} \sum_{i=1}^{n} \text{Var}(X_i) \\
& = \lim_{n \to \infty} \frac{1}{n^2} \sum_{i=1}^{n} p_i(1-p_i) \\
& \leq \lim_{n \to \infty} \frac{1}{n^2} \sum_{i=1}^{n} \frac{1}{4} \\ 
&= \lim_{n \to \infty} \frac{1}{4n} \\
& = 0
\end{align*}
Thus, by the sufficient condition of the weak law of large numbers, we have $\bar{X} = \frac{1}{n}\sum_{k=1}^{n} X_k$ converges in probability to $E(X_k) = 0$.


\end{document}